\documentclass[UTF8,zihao=-4]{ctexart}

\usepackage[a4paper,top=2.4cm,bottom=2.4cm,left=2.25cm,right=2.25cm]{geometry}
\usepackage{amsmath,amssymb}
\usepackage{booktabs,longtable,array,tabularx}
\usepackage{listings}
\usepackage{enumitem}
\usepackage{xcolor}
\usepackage{hyperref}
\usepackage{fancyhdr}

\definecolor{inkblue}{RGB}{34,76,112}
\definecolor{softgray}{RGB}{246,248,250}
\lstset{basicstyle=\small\ttfamily,breaklines=true,showstringspaces=false,frame=single,rulecolor=\color{softgray}}
\newcommand{\slot}[1]{\textcolor{inkblue}{\textit{【#1】}}}
\newcommand{\guidehead}[1]{\par\medskip\noindent\textbf{\color{inkblue}#1}\par}
\hypersetup{
  colorlinks=true,
  linkcolor=inkblue,
  urlcolor=inkblue,
  citecolor=inkblue,
  pdfauthor={数学建模学习小组},
  pdftitle={数学建模阶段学习笔记与论文写作模板}
}

\pagestyle{fancy}
\fancyhf{}
\fancyhead[L]{数学建模阶段学习笔记与论文写作模板}
\fancyhead[R]{\thepage}
\setlength{\headheight}{15pt}
\setlength{\parindent}{2em}
\setlength{\parskip}{0.25em}
\setlist{nosep,leftmargin=2.2em}

\title{\bfseries 数学建模阶段学习笔记与论文写作模板}
\author{队长：夏同、组员1：温启韬、组员2：梁文曦}
\date{2026年8月11日}

\begin{document}

\maketitle

\begin{abstract}
这份报告记下了本队现阶段的训练情况，也是刘小兰教授在队长群中布置、由丁老师转发的学习作业。

前半部分顺着思维导图和三人分工，把十三个算法与模型学习单元过了一遍。后半部分精读2024年A题一等奖论文A053《基于目标优化的板凳龙形态调节》，从摘要、假设、建模、求解一路读到结果解释和附录，看作者怎样把五个分问接成一篇文章。

\end{abstract}

\tableofcontents
\newpage

\section{阶段学习笔记}

\subsection{选题侧重与成员分工}

近六年的A、B、C题没有固定解法，训练重点却分得很清楚。A题常从物理过程推到方程，还要留意离散误差；B题多半围着变量、约束和可行方案转；C题先认字段和样本，再谈预测、评价或决策。分工就按这个差别来：夏同抓机理与连续模型，温启韬抓规划和优化，梁文曦抓数据处理与统计。我们的标准是能用自己的话说出为什么选它，能把公式还原成程序步骤，也知道它在什么条件下会失灵。

\subsection{十三个算法与模型学习单元}

\setlength{\tabcolsep}{2.5pt}
\small
\begin{longtable}{p{1.25cm}p{2.55cm}p{6.25cm}p{5.45cm}}
\caption{成员学习内容、代码逻辑与使用边界}\label{tab:learning-units}\\
\toprule
成员 & 算法或模型 & 基本原理、步骤与代码逻辑 & 适用问题、优缺点和限制\\
\midrule
\endfirsthead
\toprule
成员 & 算法或模型 & 基本原理、步骤与代码逻辑 & 适用问题、优缺点和限制\\
\midrule
\endhead
\bottomrule
\endfoot

夏同 & ODE与RK4
& 先从受力、守恒或状态转移关系写出 $\dot{y}=f(t,y)$，再由初值出发，用四个斜率的加权值向前推进。
& 适合状态随时间连续变化的力学、传热、种群和传播过程。RK4在精度与实现难度之间较为均衡，但刚性方程会受稳定域限制；此时应改用刚性求解器，不靠缩小一步掩盖问题。\\

夏同 & PDE与有限差分
& 空间区域、初始条件和边界条件要先落到网格上，再将时间、空间导数换成差分格式。
& 适合规则区域上的热、波和扩散问题。公式与网格位置对应直观，代价是复杂边界难处理，显式格式还须核对稳定条件。\\

夏同 & 有限元模型
& 从PDE的形式入手：划分单元，计算局部矩阵，装配全局稀疏矩阵，“施加”边界后求节点量。
& 适合不规则几何、变系数介质和混合边界。几何适应性强，但网格质量、形函数和边界
装配都容易引入错误；规则区域的入门题先用有限差分作基线。\\

温启韬 & LP与MILP
& 先逐项列出决策变量、线性目标、等式和不等式；遇到启停、选择或分配，才引入0-1变量。
& 适合资源分配、排程、选址和生产决策。线性模型便于解释并能给最优性信息；整数规模大时计算量上升，Big--M须来自实际上界，不能任意指定过大的常数。\\

温启韬 & 粒子群优化
& 每个粒子都保留位置、速度、个体最好解和群体最好解；迭代时包含速度更新、越界处理和重新评价。
& 适合连续、非光滑或难求导的参数搜索，编码简短；它没有连续全局最优证书，结果会受边界、罚项和随机种子影响。若变量只有一维且单调，可改用遍历或二分。\\

温启韬 & 图论与网络优化
& 建模时把对象落成节点、关系落成带权边。非负单源最短路用Dijkstra；容量与费用同时出现时用最小费用流。实战时，先建图，再运行算法，最后把路径或流量映射回现实对象并核对容量。
& 适合运输、路径、匹配和网络资源配置。结构清楚、结果易复算；动态图、负权边和多目标路线要另行处理，最小生成树也不能代替两点间最短路。\\

梁文曦 & 缺失、异常与重复处理
& 拿到附件后，先建数据字典和主键，并把未检出、未知、未采样分开记录；随后才按成因决定删除、插值或KNN。
& 它决定后续样本是否可信，不能在划分训练集前用全体数据估计均值或尺度。\\

梁文曦 & PCA与K--Means
& 数据标准化后，对协方差矩阵作特征分解并按解释率选取主成分；降维结果再送入K--Means，比较不同簇数的轮廓系数。
& PCA适合线性降维，K--Means适合近似球形簇。两者速度快、图形直观。\\

梁文曦 & 熵权与TOPSIS
& 指标须先正向化和无量纲化。熵权刻画样本差异，TOPSIS据此计算对象到正、负理想解的距离。
& 适合多指标评价与方案排序。计算过程清楚，但“差异大”不等于“业务上更重要”；指标
方向、相关重复和权重敏感性都要说明，不能只给最终名次。\\
\end{longtable}
\normalsize

\subsection{本阶段做得较好的地方}

三人的学习方向不同，但笔记用了同一套记法。读群内论文的方式是顺着一项结果回查所用公式和程序，把容易在复现时卡住的地方记下来。

\subsection{目前不足}

能把算法原理讲清楚，不等于能在竞赛时间内独立处理数据、调通程序、解释异常输出，后续训练要看实际运行记录和真题结果。

\section{一等奖论文学习分析}

\subsection{选文与阅读范围}

本次选择2024年A题《基于目标优化的板凳龙形态调节》。

\subsection{论文结构及各部分写法}

\small
\begin{longtable}{p{2.25cm}p{4.0cm}p{5.1cm}p{4.2cm}}
\caption{A053章节功能与我们理解}\label{tab:a053-structure}\\
\toprule
部分 & 论文在写什么 & 具体写法 & 我们吸收的做法\\
\midrule
\endfirsthead
\toprule
部分 & 论文在写什么 & 具体写法 & 我们吸收的做法\\
\midrule
\endhead
\bottomrule
\endfoot
摘要（p.1）
& 依次回答运动状态、碰撞、最小螺距、调头曲线和速度上限。
& 每个分问都给对象、模型动作和交付值，如第一次碰撞时刻、螺距和龙头速度；算法名放在它所求的量之后。
& 我们写摘要时也按分问落笔，但每问至少留下一个能作判断的结果；顺序词能省则省，不把摘要写成目录。\\

问题分析（p.3--4）
& 把板凳龙的整体运动拆成相邻把手位置递推，再说明各问增加的约束。
& 问题四被拆成“求较短调头曲线”和“求调头过程中的运动状态”两个子任务，前者输出路径，后者沿用位置递推。
& 我们把输入输出关系画清后，再决定方法；沿用前问时，直接列明哪些量继承、哪些量新加。\\

模型假设（p.2）
& 限定人为扰动、板凳刚体和开始调头的时机。
& 假设会进入方程或事件判定，没有写与计算无关的常识句。
& 每条假设都要回答“它使哪一项可以怎样计算”，并在评价中说明放宽后要改哪个模块。\\

符号说明（p.3）
& 给出螺距、极角、位置、速度、圆弧半径等贯穿五问的量。
& 统一公共符号，分问新增量在首次使用处定义，减少同一字母承担两个含义。
& 符号表只放多次使用的量；仅在一处使用的中间量随公式解释。\\

模型建立（p.6--21、p.23--36）
& 写螺线弧长、定长递推、矩形碰撞、相切圆弧和多段路径状态。
& 公式按“已知把手--定长方程--候选根--运动方向筛选--下一把手”推进；碰撞从几何对象落到矩形投影区间。
& 公式尽量按程序实际执行的次序出现。多根问题交代筛选条件，分段路径则先列全状态及切换条件。\\

模型求解（p.8--16、p.22--31、p.34--41）
& 用二分、时间离散、变步长遍历、状态递推和粒子群得到五问结果。
& 螺距先在 $[45,46]\,\mathrm{cm}$ 内定位，再以 $0.01\,\mathrm{cm}$ 步长细查；碰撞检测在第一次触发时停止。
& 粗查负责圈定“边界在哪一段”，细查再回答“边界精确到多少”；随机算法另报种子和重复次数。\\

结果及分析（p.22、p.27--28、p.32--33、p.38、p.41）
& 报告碰撞时刻、螺距、圆弧半径、曲线长度和速度上限，并解释局部异常。
& 参数微变后若第一次碰撞对象改变，论文从局部几何解释输出跳变；12--13秒的速度突增则用相同时间内的路程差解释。
& 表后的正文要回答“为何变、哪条约束起作用、结论只到什么范围”。\\

检验与边界
& 检查首次碰撞对象、全过程可行性和现实动作限制。
& p.17用反设和逐节前推说明最早碰撞对象；p.28用“终点可行但中途碰撞”的方案否定只查终点；p.32从龙身不能倒退推出半径下界。
& 检验要对准一个具体疑问。回代、反例和物理边界各有侧重。\\

附录（p.43--62）
& 给出逐把手循环、二分、分段路径、碰撞和粒子群代码。
& 正文四种运动状态与代码状态码1--4相互对应，表格输出也由固定变量名生成。
& 附录须留下复现正文结果所需的入口、参数和输出说明；正文讲关键流程，不整段粘贴程序。\\
\end{longtable}
\normalsize

\subsection{五个分问的思路怎样连起来}

\begin{enumerate}
  \item 问题一先确定龙头在阿基米德螺线上的位置，相邻把手随后按距离不变的关系逐节递推；
  速度由短时间内的弧长变化得到。除了题目要求的数值，这一问还写出了后四问都会调用的位置和速度计算方法。
  \item 问题二把每块板凳写成矩形，用分离轴判断重叠。检查对象并非凭经验删减，
  而是先用反证法确定哪一类构件最早碰撞，再让程序在首次触发处停止。
  \item 问题三中的“能进入调头区域”指整个盘入过程均不碰撞。搜索先圈出螺距所在区间，再改用较小步长。
  有些候选在终点看似可行，中途却已发生碰撞，因此判断不能只看最后一帧。
  \item 问题四先由相切关系写出两段圆弧。若半径过小，后续构件就得倒退让位，由这一实际限制可得大圆弧半径的下界。
  路径确定后，位置递推按盘入、两段圆弧和盘出四种状态继续运行。
  \item 问题五沿用问题四的运动模型，把各把手速度改写为对龙头速度的限制；粒子群在这里只搜索允许的龙头速度，
  板凳龙的运动仍按问题四计算。
\end{enumerate}

\subsection{最值得模仿的段落组织}

\paragraph{从局部关系写到全局运动。}
板凳龙的构件很多，反复用到的却是相邻把手间的固定距离。求出一个把手后，下一把手的位置由定长方程给出；若方程有多个根，还要按构件次序和运动方向逐个排除。这种写法紧跟实际计算顺序，读到公式时也能看出程序下一步如何推进。

\paragraph{先证明检查范围不会漏，再缩小程序。}
讨论第一次碰撞时，论文先假设第 $i$ 个构件更早碰撞，再利用各构件经过同一轨迹的先后关系逐节前推，最后同龙头的运动状态相矛盾。于是，程序只检查与龙头有关的候选对象便有了依据。我们以后缩小搜索或检查范围时，也要先回答会不会漏解，计算量的减少放在后面说明。

\paragraph{用反例判断并简化方案。}
螺距约为 $42\,\mathrm{cm}$ 的方案到达终点时看似可行，盘入途中却已经碰撞。这个反例说明终点状态不能代替全过程检查。论文处理12--13秒的速度突增时，则先找到节点所在的小圆弧，再比较相同时间内相邻节点走过的路程。异常时段、局部位置和路程差依次对应起来以后，速度变化才有了几何解释。


\clearpage
\section{本队论文模板}

这里的“模板”包括标题树、写前账本和核对动作，不提供一套可以逐句替换名词的成品正文。训练时先用表格查缺漏，动笔后再按本题的计算顺序取舍材料；没有触发的账本不进入正文，已经写清的关系也不因模板存在而重复一遍。可直接编辑的双模式文件位于\texttt{thesis\ template/main.tex}，本节保留的是课程作业所需的使用说明和填写接口。

\subsection{按竞赛规章要求的提交版式}

\begin{longtable}{p{2.7cm}p{5.7cm}p{6.2cm}}
\caption{2026 年论文与支撑材料硬性要求}\label{tab:official-format-rules}\\
\toprule
对象 & 纸质版/正文要求 & 电子版/支撑材料要求\\
\midrule
\endfirsthead
\toprule
对象 & 纸质版/正文要求 & 电子版/支撑材料要求\\
\midrule
\endhead
纸张和装订 & 白色 A4 纸，单面或双面均可；四边页边距至少 $2.5\,\mathrm{cm}$，左侧装订。 & 电子论文必须与纸质版内容及格式（含附录）一致。\\
首页顺序 & 第 1 页承诺书，第 2 页编号专用页，第 3 页摘要专用页。 & 电子论文不放承诺书和编号专用页，第一页直接是摘要专用页。\\
摘要与页码 & 摘要含标题、关键词，无需英文，原则上不超过 1 页；从摘要页起以阿拉伯数字从 1 连续编号，页码居页脚中部。 & 参赛论文为一个单独 PDF 或 Word 文件，建议 PDF，文件不超过 20 MB，不能再压缩后上传。\\
正文与附录 & 正文从纸质版第 4 页开始，不设目录，不超过 30 页；正文后为附录，页数不限，须打印并与正文装订。 & 支撑材料另交一个 RAR 或 ZIP 文件，不超过 20 MB。\\
附录和程序 & 附录列出支撑材料文件，并纳入全部完整、可运行的源程序（含 Excel、SPSS 等软件的交互命令）；确实未用程序时写明“本论文没有用到程序”。 & 支撑材料至少包括全部可运行源程序、自主查阅数据（赛题原始数据除外）和较大篇幅中间图表；确实没有时在附录注明“本论文没有支撑材料”。\\
匿名与引用 & 摘要、正文和附录不得出现参赛者身份、学校或赛区信息；他人及公开资料须在正文标注并列入参考文献。 & 支撑材料同样不得出现身份、学校或赛区信息，也不得放入承诺书和编号专用页。\\
\bottomrule
\end{longtable}

竞赛期间必须由本队队员独立完成赛题。指导教师不能解释赛题、提供选题或解题建议、提供资料、修改论文或给出修改意见；队员也不能同队外任何人讨论赛题，不能在社交平台、论坛、代码平台或直播间浏览、发布、讨论当届赛题内容。AI 工具可以作为辅助，但作品的原创性、真实性和准确性仍由参赛队负责，并须遵守专门披露规定。

\subsubsection{AI 工具披露位置与材料}

2026 年试行规定要求参赛队公开、如实说明 AI 工具的使用情况，并逐项人工审查和核实相关内容。声明不是附录，也不放在参考文献之后，而是作为独立部分排在参考文献之前；两种情况只能选择其一。

\begin{longtable}{p{2.7cm}p{5.7cm}p{6.2cm}}
\caption{AI 工具使用披露清单}\label{tab:official-ai-rules}\\
\toprule
实际情况 & 论文内处理 & 支撑材料与人工责任\\
\midrule
\endfirsthead
\toprule
实际情况 & 论文内处理 & 支撑材料与人工责任\\
\midrule
\endhead
未使用任何 AI 工具 & 在参考文献之前设置“AI 工具使用声明”，使用规定给出的未使用声明。 & 不提交虚构记录；声明必须与竞赛期间的实际过程一致。\\
使用 AI 工具 & 在参考文献之前设置“AI 工具使用声明”，简要写明用途，并指向支撑材料。 & 支撑材料必须包含独立 PDF，文件名为\texttt{AI 工具使用详情.pdf}；写明工具名称及版本或型号、使用目的和环节、主要提示方式与使用过程、采纳情况以及人工修改与核验。语言润色可不写最后一项。\\
核心建模与分析 & 核心建模与分析必须由参赛队主导；AI 参与完成的内容须逐项人工审查与核实。 & 故意隐瞒、虚假声明，或把未经必要人工审查和核实的 AI 生成内容直接作为核心成果提交，将被取消评奖资格。\\
\bottomrule
\end{longtable}

正式论文在参考文献前保留下面这一段，按实际情况删去另一项：
\begin{quote}\small
\textbf{AI 工具使用声明}

未使用时：本参赛队在竞赛过程中未使用任何 AI 工具。

使用时：本参赛队在竞赛过程中使用了 AI 工具，主要用于\slot{简要用途，如语言润色、代码调试等}，详细使用情况见支撑材料。
\end{quote}
规定没有要求把 AI 工具另列为参考文献，也没有要求在正文中逐句标出 AI 辅助痕迹；但是支撑材料中的用途、过程和人工核验必须能够与论文及代码相互对应。该规定自 2026 年 9 月 1 日起试行，此前规定与之不一致时以本规定为准。

\subsection{论文中的目录与篇幅}

下表来自本次核读的 59 篇 A/B/C 题编号论文。表中数字是章节标题命中后的页级跨度统计，用来判断一部分通常写到什么程度，不是要求各章照着中位数凑页数；检出率较低的检验、灵敏度和改进部分只作定性参考。

\begin{longtable}{p{2.3cm}p{1.5cm}p{3.0cm}p{8.0cm}}
\caption{59 篇 A/B/C 论文的章节跨度与交付物}\label{tab:template-corpus-length}\\
\toprule
章节 & 检出数 & 页级跨度中位数 [Q1,Q3] & 正式稿中可核对的内容\\
\midrule
\endfirsthead
\toprule
章节 & 检出数 & 页级跨度中位数 [Q1,Q3] & 正式稿中可核对的内容\\
\midrule
\endhead
摘要与关键词 & 58/59 & 1 [1,1] & 分问任务、实际做法、关键数值、检查方法、结论和 3--5 个关键词。\\
问题重述 & 57/59 & 1 [1,1] & 研究对象、输入、输出、硬约束，以及分问之间传递的量。\\
问题分析 & 57/59 & 1 [1,1] & 从题面、数据或物理边界出发的取舍依据，以及为何保留或放弃某条路线。\\
模型假设 & 48/59 & 1 [1,1] & 每条假设的对象、条件、计算后果和适用边界。\\
符号说明 & 58/59 & 1 [1,3.75] & 变量、参数、集合、下标、单位；只放跨公式复用的量。\\
模型建立 & 58/59 & 2 [1,6] & 定义、方程/目标函数、约束、边界条件和与题意的映射。\\
模型求解 & 55/59 & 2 [1,6] & 初始化、更新、可行性处理、搜索范围、停止条件、输出，以及代码中的对应函数。\\
结果分析 & 47/59 & 7 [1,12] & 读数、方向、数量级、变化原因、约束状态、现实含义和结果边界。\\
模型检验 & 22/59 & 1.5 [1,7] & 检验对象、独立证据、指标/阈值、结果和判断；缺少这些内容时，不宜声称“验证有效”。\\
灵敏度分析 & 25/59 & 1 [1,2] & 基准、扰动来源、范围、固定量、输出变化和决策是否切换。\\
模型评价 & 51/59 & 1 [1,1] & 有公式或数据支撑的优点，以及受对象、数据和计算预算限制的缺点。\\
改进方案 & 19/59 & 1 [1,1] & 新增数据/关系、修改层、预计改善量和仍需检验的边界。\\
参考文献 & 57/59 & 1 [1,11] & 正文真正使用过的题面、方法、数据和软件来源。\\
附录 & 48/59 & 16 [6,28] & 代码入口、参数、数据字典、运行顺序、补充推导、完整结果和输出文件。\\
\bottomrule
\end{longtable}

\subsection{正式稿标题次序与分问承接关系}

正式稿可按下列顺序组织；“问题一”等应换成题目中的实际编号和任务名，标题通常写实际任务，无需另列算法名称。写到下一问时，应点明沿用了上一问的哪些变量、参数或筛选结果：

\begin{enumerate}
  \item 摘要与关键词；
  \item 问题重述；
  \item 问题分析；
  \item 模型假设；
  \item 符号说明；
  \item 问题一：\slot{任务名}；问题二：\slot{任务名}；……每问均包含“模型建立”与“模型求解”；
  \item 结果分析与模型检验；
  \item 灵敏度/稳健性分析（若题目或结果支持）；
  \item 模型评价与改进；
  \item AI 工具使用声明：按实际情况选择官方规定中的一种表述；
  \item 参考文献；
  \item 附录：代码、数据、补充结果和运行说明。
\end{enumerate}

\begin{longtable}{p{1.2cm}p{3.0cm}p{3.0cm}p{3.0cm}p{4.2cm}}
\caption{分问任务表：先列承接量，再写正文}\label{tab:template-task-matrix}\\
\toprule
分问 & 输入/证据 & 要求输出 & 硬约束或计算标准 & 传给下一问的量\\
\midrule
\endfirsthead
\toprule
分问 & 输入/证据 & 要求输出 & 硬约束或计算标准 & 传给下一问的量\\
\midrule
\endhead
问题一 & \slot{附件字段、初始状态、题面量} & \slot{数值/分类/曲线/方案} & \slot{单位、范围、边界} & \slot{变量、参数、筛选结果}\\
问题二 & \slot{沿用问题一的哪些输出} & \slot{……} & \slot{新增约束} & \slot{……}\\
问题三 & \slot{……} & \slot{……} & \slot{……} & \slot{最终决策或评价准则}\\
问题四（如有） & \slot{……} & \slot{……} & \slot{……} & \slot{……}\\
\bottomrule
\end{longtable}

\subsection{摘要与关键词：一页内说完做法和结论}

\guidehead{摘要要留下的信息}先点明研究对象和条件，再按分问写实际做法、最能区分方案的数值或关系、所用检查以及结论适用范围。算法名称随求解动作出现即可；摘要中的数字均应来自正文计算或程序输出。

\guidehead{写前先备齐的材料}摘要不使用统一句架。先在草稿旁列出研究对象和总限制，再为每问各留四项：实际做法、最能区分方案的输出、数字所在表格或程序文件、与该输出相称的检查。写成正文时，核心分问可以展开两三句，沿用前问的部分则压成一句；不要让每问保持相同字数和语法。关键词从题目对象、主模型和求解任务中选取 3--5 个，不把四项材料依次改写成四个并列短语。

\guidehead{核对}每个分问至少留下一项可复算的输出；数字能在正文表格或程序输出中找到；单位前后一致；检查方法的表述与实际工作相称。“最优、稳定、准确”等判断还要带上比较对象、指标和范围。

\subsection{问题重述：把题面整理成可计算的任务}

\guidehead{重述时依次核对的对象}先核对研究对象和时间、空间范围，再按题面顺序写各问的已知量与所求量。硬约束、单位和附件字段放在第一次影响任务的位置；预测值与观测值、局部量与全局量等易混口径只在确实存在时另句区分。背景很短时可直接进入问题一，不必先写一段意义介绍。

\guidehead{注意}这一节照题意说明任务，算法选择放到后面的问题分析中。对象范围、时间窗、单位以及题面中“必须”“可以”等限定词应保持原意；附件里没有的字段仍按未知量处理。

\subsection{问题分析：写出方法从哪里来}

问题分析要回答“为何这样处理”，但无需抄录杂乱的试算过程。正文可以从一个异常读数写起，也可以先指出某个边界不满足，或说明上一问留下的变量怎样进入本问。只要后面的选择能在题面、数据、公式、图表或运行记录中找到依据，段落次序便可随问题本身变化。

\begin{longtable}{p{2.0cm}p{3.5cm}p{3.4cm}p{3.3cm}p{3.4cm}}
\caption{问题分析记录表：保留实际取舍}\label{tab:template-analysis-ledger}\\
\toprule
分问 & 观察到的事实/限制 & 尚未解决的量或矛盾 & 采取的处理及其影响 & 带入模型的范围或变量\\
\midrule
\endfirsthead
\toprule
分问 & 观察到的事实/限制 & 尚未解决的量或矛盾 & 采取的处理及其影响 & 带入模型的范围或变量\\
\midrule
\endhead
问题一 & \slot{图、表、物理条件或失败试算} & \slot{要解释/求解的量} & \slot{保留、舍弃、改写或分段} & \slot{变量、范围、初值}\\
问题二 & \slot{问题一的残差/新约束} & \slot{……} & \slot{……} & \slot{……}\\
问题三 & \slot{决策冲突或不确定性} & \slot{……} & \slot{……} & \slot{……}\\
\bottomrule
\end{longtable}

记录表只在发生真实取舍时填写，正文也不必按表的五列依次复述。从触发选择的那一列起笔：可能是附件中的异常字段、题面的物理限制、一次不可行试算，也可能只是上一问留下的状态。没有运行过旧方案时，不写“原方法存在不足”；后问仅沿用已有结果时，直接点明输入和新增约束即可。

\subsection{模型假设：每条假设都要产生计算后果}

\begin{longtable}{p{3.2cm}p{3.1cm}p{4.0cm}p{3.6cm}}
\caption{模型假设记录表}\label{tab:template-assumption-ledger}\\
\toprule
假设（对象+条件） & 题面/数据依据 & 在公式或程序中的后果 & 检查与失效边界\\
\midrule
\endfirsthead
\toprule
假设（对象+条件） & 题面/数据依据 & 在公式或程序中的后果 & 检查与失效边界\\
\midrule
\endhead
\slot{忽略某项影响} & \slot{题面允许/量级比较} & \slot{删除哪一项或固定哪个量} & \slot{在何种范围失效}\\
\slot{连续/刚体/独立等条件} & \slot{观测或定义} & \slot{改变哪种状态转移或误差} & \slot{用何种回代/扰动检查}\\
\bottomrule
\end{longtable}

正式稿不把每条假设写成同一种四句说明。假设较直观时，将对象、条件和计算后果合成一句；只有量级比较、适用范围或失效后果不明显时再展开。若某条“假设”没有删去方程项、固定变量、改变状态转移或限定数据范围，它多半不必单列。

\subsection{符号说明：只收录会反复使用的量}

\begin{longtable}{p{2.2cm}p{7.0cm}p{2.2cm}p{3.3cm}}
\caption{符号说明可填写表}\label{tab:template-symbol-ledger}\\
\toprule
符号 & 含义（对象和方向） & 单位/取值域 & 首次出现位置\\
\midrule
\endfirsthead
\toprule
符号 & 含义（对象和方向） & 单位/取值域 & 首次出现位置\\
\midrule
\endhead
$x_i$ & \slot{第 $i$ 个对象的决策/观测量} & \slot{单位或域} & 式(\slot{})\\
$\theta$ & \slot{待估参数/情景参数} & \slot{单位或范围} & 式(\slot{})\\
$\mathcal{I}$ & \slot{对象集合/时间集合} & \slot{集合定义} & 式(\slot{})\\
\bottomrule
\end{longtable}

\subsection{分问题模型建立与求解：逐问填写}

下表先检查各问是否闭合，不决定正文的段落顺序。变量和方程可以连写，边界条件可以跟在首次用到它的公式之后；求解器的输入、状态和停止条件则放在程序真正接管的位置。

\begin{longtable}{p{1.8cm}p{2.4cm}p{2.7cm}p{2.7cm}p{2.8cm}p{2.1cm}}
\caption{分问模型填写表}\label{tab:template-model-card}\\
\toprule
分问 & 变量/状态 & 目标或方程 & 约束/边界 & 求解动作 & 输出及核查\\
\midrule
\endfirsthead
\toprule
分问 & 变量/状态 & 目标或方程 & 约束/边界 & 求解动作 & 输出及核查\\
\midrule
\endhead
问题一 & \slot{} & \slot{} & \slot{} & \slot{初始化--更新--停止} & \slot{}\\
问题二 & \slot{} & \slot{} & \slot{} & \slot{} & \slot{}\\
问题三 & \slot{} & \slot{} & \slot{} & \slot{} & \slot{}\\
\bottomrule
\end{longtable}

\guidehead{每一问可调用的材料块}下面六块用于检查，不是六个固定段落，也不要求按此顺序出现。
\begin{enumerate}
  \item 题面或前问给了哪些量，其中哪些能直接代入，真正需要求的量还剩什么；
  \item 首次出现的变量、取值域和单位，以及不参与决策但在本问保持不变的量；
  \item 守恒、几何、统计定义或业务规则怎样落到方程、目标和约束；
  \item 多根、分支、越界、不可行解或缺失字段怎样裁决；没有这些情况时不补写；
  \item 程序从哪里开始，怎样更新，在什么条件下停止，结果写到哪个文件、表或图；
  \item 本问实际检查过的误差、边界或失败状态。没有额外证据时，不另设一段通用“模型边界”。
\end{enumerate}
正文标题写本问的实际任务。若推导自然地从已知量进入公式，可直接定义变量并列式；若一次试算改变了搜索范围，就把取舍留在试算之后。代码入口、函数名、种子和完整参数表放在求解段末尾或附录，主文只保留理解结果所需的部分。

\subsection{A/B/C 题型}

\begin{longtable}{p{1.6cm}p{4.4cm}p{4.5cm}p{4.3cm}}
\caption{三类题型常需补充的内容}\label{tab:template-abc-slots}\\
\toprule
题型 & 建模时需说明 & 求解时需说明 & 结果与检验需说明\\
\midrule
\endfirsthead
\toprule
题型 & 建模时需说明 & 求解时需说明 & 结果与检验需说明\\
\midrule
\endhead
A：机理/连续 & 状态、坐标、初值、边界、守恒或受力关系；离散格式与连续模型分开。 & 步长/网格、稳定条件、初值推进、事件触发、收敛或误差界；多根按物理顺序筛选。 & 回代、守恒、边界和量纲；异常时段回到局部轨迹或同时间路程，一条结果曲线不足以解释原因。\\
B：综合优化 & 决策变量、目标层级、硬约束、可行域和前后问承接量；多目标的优先级或权重来源。 & 编码、初始解、越界/不可行处理、搜索范围、停止条件、重复协议和一致的比较条件。 & 目标值、约束残差、候选位置、成本与现实含义；改善幅度较小时同时报告最大值和平均值。\\
C：数据决策 & 原字段--清洗/聚合--特征--标签--划分--输出关系；指标方向、业务语义和缺失类型。 & 外层测试保持独立，预处理只读训练折；说明模型损失、阈值、类别不平衡、随机重复，以及预测结果怎样进入决策。 & 区分训练、验证和测试数据；报告基准方法、混淆矩阵或误差、风险/收益和场景复评。同源指标不能替代外部验证，只能反映内部一致性。\\
\bottomrule
\end{longtable}

\guidehead{三类题的自然起点}A 类常从边界、受力、守恒、多根或首次事件进入；B 类可从一条最先被破坏的硬约束、状态规模或成本冲突进入；C 类则多从样本单位、字段语义、分布形态或训练测试边界进入。表~\ref{tab:template-abc-slots} 中的内容按当前题目选取，不把 A/B/C 三行各改写成一段，也不要求每问都同时写建模、求解、结果和检验四层说明。


\subsection{结果分析与模型检验：表格之后还要解释}

\begin{longtable}{p{2.2cm}p{3.0cm}p{3.0cm}p{3.0cm}p{3.5cm}}
\caption{结果解释记录表}\label{tab:template-result-ledger}\\
\toprule
结果对象 & 条件/对象 & 读数与单位 & 与谁比较 & 原因、影响和边界\\
\midrule
\endfirsthead
\toprule
结果对象 & 条件/对象 & 读数与单位 & 与谁比较 & 原因、影响和边界\\
\midrule
\endhead
\slot{表/图/输出编号} & \slot{情景、样本、时段} & \slot{数值、方向、数量级} & \slot{基准方法/候选/实测} & \slot{原因、约束、现实动作和适用范围}\\
\slot{} & \slot{} & \slot{} & \slot{} & \slot{}\\
\bottomrule
\end{longtable}

一张表或图不必配一段完整的五步说明。先圈出会改变判断的行、时段或对象：普通读数一句交代，异常、拐点、约束活跃或方案切换再展开原因。表中已经显然可见的大小关系不逐格复述；只有实际检查到适用边界时才写边界句。

\begin{longtable}{p{2.6cm}p{3.4cm}p{3.0cm}p{3.0cm}p{2.8cm}}
\caption{模型检验记录表}\label{tab:template-validation-ledger}\\
\toprule
检验问题 & 独立证据/数据 & 指标与接受标准 & 结果 & 裁决及影响范围\\
\midrule
\endfirsthead
\toprule
检验问题 & 独立证据/数据 & 指标与接受标准 & 结果 & 裁决及影响范围\\
\midrule
\endhead
\slot{守恒/回代是否成立} & \slot{未参与拟合的量/公式} & \slot{误差定义、阈值} & \slot{数值} & \slot{保留/修正哪一层}\\
\slot{预测/分类是否可用} & \slot{独立测试或交叉验证} & \slot{MAE/RMSE/召回等} & \slot{} & \slot{}\\
\slot{方案是否可执行} & \slot{约束复算/场景} & \slot{违约率、成本、风险} & \slot{} & \slot{}\\
\bottomrule
\end{longtable}

检验实施前先在记录表中确定指标、阈值、满足条件时保留的部分，以及不满足时需要回到哪一层修改；正文根据实际发生的裁决来写，不照抄一条统一的条件句。训练回代、内层调参、最终独立测试是三件事，得分开。一次重复运行只能说明这一次的结果，一条相关曲线也撑不起因果结论。

\subsection{灵敏度与稳健性：改变什么、为什么改变、何时换方案}

\begin{longtable}{p{2.4cm}p{3.0cm}p{3.0cm}p{3.0cm}p{3.5cm}}
\caption{灵敏度/稳健性记录表}\label{tab:template-sensitivity-ledger}\\
\toprule
扰动量 & 基准与来源 & 扰动范围/固定量 & 输出变化 & 决策切换与结论边界\\
\midrule
\endfirsthead
\toprule
扰动量 & 基准与来源 & 扰动范围/固定量 & 输出变化 & 决策切换与结论边界\\
\midrule
\endhead
\slot{参数/误差/场景} & \slot{测量、估计或题面来源} & \slot{范围、步长、其余固定} & \slot{最大/平均/排序/曲线} & \slot{是否换方案，阈值在哪里}\\
\slot{} & \slot{} & \slot{} & \slot{} & \slot{}\\
\bottomrule
\end{longtable}

正文按实际结果选择重点：连续变化主要报告方向和幅度，存在阈值时报告切换位置，排序未变时说明测试范围即可。基准、扰动来源和固定量可合并在图题、表头或段首，不必每次铺成相同四句；没有观察到切换时不要补造一个“超过阈值后”的分支。

\subsection{模型评价与改进：说明优缺点出现在哪里}

\begin{longtable}{p{2.3cm}p{4.0cm}p{3.6cm}p{3.8cm}}
\caption{模型评价与改进记录表}\label{tab:template-evaluation-ledger}\\
\toprule
评价方面 & 已有证据 & 具体优点或不足 & 可执行的改进及新检验\\
\midrule
\endfirsthead
\toprule
评价方面 & 已有证据 & 具体优点或不足 & 可执行的改进及新检验\\
\midrule
\endhead
精度/解释 & \slot{误差、回代或图表} & \slot{对哪些对象有效/在哪里失效} & \slot{新增数据、关系或校准方法}\\
计算/复现 & \slot{复杂度、运行记录、代码} & \slot{瓶颈和依赖} & \slot{缓存、分层求解或更换离散格式}\\
现实性/外推 & \slot{约束和场景复评} & \slot{未观测对象、时窗或风险} & \slot{追加样本、边界场景和接受阈值}\\
\bottomrule
\end{longtable}

“速度快、精度高、适用范围广”单独列出来并不能说明优点在哪里。评价时从表中选有运行记录或误差证据的项目，不为完整而凑齐三类；改进则点明要增加的数据或关系、受影响的模型层和复查指标。现有数据没有支撑的内容仍属于后续工作，应与本文已经得到的结果分开。

\subsection{参考文献与附录：保留复算所需信息}

正文首次使用外部定义、算法、数据或软件时标出来源，文后不列未使用的条目。附录按实际运行的“输入--处理--求解--输出”顺序排列，并填写下面各项：

\begin{longtable}{p{3.0cm}p{4.2cm}p{4.2cm}p{3.2cm}}
\caption{附录复现清单}\label{tab:template-appendix-ledger}\\
\toprule
项目 & 需给出的内容 & 与正文的对应关系 & 核对状态\\
\midrule
\endfirsthead
\toprule
项目 & 需给出的内容 & 与正文的对应关系 & 核对状态\\
\midrule
\endhead
数据字典 & 文件名、字段、单位、缺失/异常编码、主键 & 问题重述、预处理和符号表 & \slot{待核/已核}\\
代码入口 & 主脚本、函数调用顺序、依赖和版本 & 模型求解段、输出表/图 & \slot{}\\
参数与随机性 & 初值、范围、步长、种子、重复次数、停止条件 & 模型卡和灵敏度表 & \slot{}\\
中间结果 & 关键状态、候选、约束残差、日志或完整结果表位置 & 结果分析和检验 & \slot{}\\
最终文件 & 运行命令、输出路径、文件哈希或版本日期 & 摘要数字和交稿文件 & \slot{}\\
\bottomrule
\end{longtable}




\end{document}
